\documentclass[border=0pt,varwidth=false]{standalone}
\usepackage{ews-figures}
\begin{document}
\begin{ewsfigure}

\ewsband{A}{The two implemented architectures}

\begin{ewsdiagram}
\begin{tikzpicture}[
  every node/.style={font=\fontsize{7.2}{8.6}\selectfont},
  role/.style={draw=ewsrule, fill=white, rounded corners=1.5pt,
               inner sep=3pt, minimum height=14pt, minimum width=46pt,
               align=center},
  store/.style={role, draw=ewsblue, fill=ewsblue!6},
  shared/.style={role, draw=ewsaccent, fill=ewsaccent!6,
                 minimum width=158pt},
  coord/.style={role, draw=ewsslate, fill=ewsslate!6, minimum width=52pt},
  flow/.style={-{Stealth[length=4pt]}, draw=ewsmute, line width=0.5pt},
  bi/.style={{Stealth[length=4pt]}-{Stealth[length=4pt]},
             draw=ewsmute, line width=0.5pt}]

% ---------------- Panel A: one shared transcript
\node[anchor=west, font=\ewslabel\fontsize{7.4}{9}\selectfont]
  (la) at (0,0) {PANEL A \textemdash{} SHARED TRANSCRIPT};
\node[role, below=9pt of la.west, anchor=north west] (sup) {Supervisor};
\node[role, right=7pt of sup] (jus) {Justine};
\node[role, right=7pt of jus] (jul) {Juliette};
\node[shared, below=11pt of jus] (tr) {one accumulated transcript\\
  every role reads and writes the same state};
\draw[bi] (sup.south) -- ($(sup.south|-tr.north)$);
\draw[bi] (jus.south) -- (tr.north);
\draw[bi] (jul.south) -- ($(jul.south|-tr.north)$);

% ---------------- Panel B: separated memories
\node[anchor=west, font=\ewslabel\fontsize{7.4}{9}\selectfont]
  (lb) at (248pt,0) {PANEL B \textemdash{} SEPARATED MEMORIES};
% Roles are spaced to clear the two-line memory boxes beside them.
\node[role, below=9pt of lb.west, anchor=north west] (sup2) {Supervisor};
\node[role, below=8pt of sup2] (jus2) {Justine};
\node[role, below=8pt of jus2] (jul2) {Juliette};
\node[store, right=10pt of sup2] (m1) {private\\memory};
\node[store, right=10pt of jus2] (m2) {private\\memory};
\node[store, right=10pt of jul2] (m3) {private\\memory};
\node[coord, left=13pt of jus2] (co) {Coordinator};
\foreach \a/\b in {sup2/m1, jus2/m2, jul2/m3} \draw[bi] (\a) -- (\b);
\foreach \t in {sup2, jus2, jul2} \draw[flow] (co.east) -- (\t.west);
\end{tikzpicture}

\ewsdiagramnote{Panel~B is genuine application-layer memory separation, not
several role names inside one context. The two harnesses do differ beyond
topology, however: prompts, demonstrations, message content and reveal order
change as well, so they do not isolate topology as a cause.}
\end{ewsdiagram}

\vspace{5pt}
\ewsband{B}{What the model produced}
\input{figures/excerpts/f4}

\vspace{5pt}
\ewsband{C}{What this establishes, and what it does not}

\begin{tcolorbox}[ewscallout]
\fontsize{8}{10.2}\selectfont
\textbf{Establishes.} Two working multi-agent harnesses, one shared-state and
one with genuinely separated memories, and role- and pressure-conditioned
safety-relevant behaviour recorded in both. \emph{The submitted trace shown
above records the refusal pole}: under coercive self-preservation framing the
model declines to produce a false admission. The compliance pole appears in the
notebook and in the screenshots reproduced in Appendix~D.
\par\vspace{3pt}
\textbf{Does not establish.} That topology causes alignment collapse, or that
the coerced-confession behaviour is typical. The notebook labels the broad
failure family 100\% reproducible while explicitly qualifying the
coerced-confession subtype as ``quite rare and less persistent than other
vectors''; that 100\% label is a breadth marker, not a confession rate.
\end{tcolorbox}

\end{ewsfigure}
\end{document}
